\title{Group Sequence Policy Optimization}

\begin{abstract}
We present Group Sequence Policy Optimization (GSPO), a reinforcement learning method designed for training large language models with enhanced stability, efficiency, and effectiveness. Departing from earlier approaches that employ token-level importance ratios, GSPO instead leverages sequence-level importance ratios to govern sequence-level processes such as clipping, reward allocation, and optimization. Our empirical analysis reveals that GSPO surpasses the GRPO algorithm both in training performance and efficiency, particularly promoting stability during Mixture-of-Experts (MoE) reinforcement learning. Additionally, GSPO streamlines the development of RL systems. These advantages have underpinned the notable advancements observed in the most recent Qwen3 models.
\end{abstract}

\section{Introduction}
\label{sec:intro}

Reinforcement learning (RL) has become a crucial framework for advancing the scalability of language models \citep{o1,dpsk-r1,qwq32b,qwen3}. By leveraging large-scale RL, language models acquire the proficiency to address complex tasks—including competitive mathematics and programming—through extended and more intricate reasoning.

For RL to be effectively scaled with increasing computational resources, it is essential to ensure consistent and robust training behavior. Nonetheless, the leading RL methodologies, such as GRPO \citep{grpo}, encounter pronounced stability challenges when applied to extremely large language models. These issues frequently lead to severe and sometimes irreversible degradation of the model, also known as model collapse \citep{qwen3, minimax-m1}. Such instability poses a significant obstacle to further enhancing language model capabilities via sustained RL-based training.

In this work, we observe that the instability present in GRPO arises from a fundamental misuse and subsequent invalidation of importance sampling weights within its algorithmic framework. Such design issues lead to substantial high-variance noise during training, which accumulates as the generated response grows longer, and is exacerbated by the clipping mechanism, eventually causing the model to collapse.

To overcome these principal shortcomings, we introduce \textbf{Group Sequence Policy Optimization (GSPO)}, a novel reinforcement learning algorithm tailored for the training of large language models. The central contribution of GSPO is its theoretically sound formulation of the importance ratio, grounded in sequence likelihood \citep{click} and consistent with the foundations of importance sampling. Moreover, GSPO determines normalized rewards as advantages across multiple responses for each query, thereby fostering coherence between sequence-level reward assignment and the optimization objective.

Our experimental results highlight the pronounced advantages of GSPO compared to GRPO, notably in terms of training stability, efficiency, and overall effectiveness. Importantly, GSPO inherently overcomes the stability problems commonly encountered during RL training of large Mixture-of-Experts (MoE) models, thus removing the dependence on intricate stabilization techniques and suggesting opportunities to streamline RL system architecture. These strengths of GSPO have directly supported the notable performance gains observed in the most recent iterations of the Qwen3 models. Looking forward, we anticipate that GSPO will serve as a robust and scalable core methodology, fostering further progress in large-scale RL training with language models.

\section{Preliminaries}

\paragraph{Notation}
Throughout this work, we represent an autoregressive language model with parameters $\theta$ as a policy $\pi_\theta$. We let $x$ signify a query and designate $\mathcal{D}$ as the collection of such queries. For any generated response $y$ corresponding to the query $x$, its conditional probability according to $\pi_\theta$ is expressed as $\pi_\theta (y | x)=\prod_{t=1}^{|y|} \pi_\theta (y_t | x, y_{<t} )$, where $|y|$ is the total token count of $y$. Additionally, we define a scoring function $r$, referred to as the verifier, which assigns a reward $r(x, y) \in [0, 1]$ to each query-response pair $(x, y)$.

\paragraph{Proximal Policy Optimization (PPO)} Using samples generated from the old policy $\pi_{\theta_\text{old}}$, PPO \citep{ppo} constrains the policy update within a proximal region of the old policy through the clipping mechanism. Specifically, PPO employs the following objective for policy optimization (we omit the KL regularization term hereinafter for brevity, as it is not the focus of this paper):

\begin{align}
\mathcal{J}_\text{PPO}(\theta) = \mathbb{E}_{ x \sim \mathcal{D},\, y \sim \pi_{\theta_\text{old}}( \cdot | x) }
\left[ \frac{1}{|y|} \sum_{t=1}^{|y|} 
\min \left( w_{t}(\theta) \widehat{A}_{t},  \, \mathrm{clip} \left( w_{t}(\theta), 1 - {\varepsilon}, 1 + {\varepsilon}\right) \widehat{A}_{t} \right)
\right],
\end{align}

Here, the importance ratio for the token $y_t$ is specified as $w_{t}(\theta) = \frac{ \pi_{\theta} (y_{t} | x, y_{<t}) }{ \pi_{\theta_\text{old}} (y_{t} | x,y_{<t})}$, the advantage $\widehat{A}_{t}$ associated with $y_t$ is computed using a separate value model, and $\varepsilon$ serves as the clipping threshold for importance ratios.

A significant obstacle for PPO in practical applications is its strong dependence on the value model. The value model is typically comparable in scale to the policy model, which results in notable demands on computational and memory resources. In addition, the success of the algorithm is closely tied to the quality of its value prediction. Achieving an accurate value model is already a complex endeavor, and the difficulty is further compounded when adapting it to handle longer outputs or more sophisticated tasks.

\paragraph{Group Relative Policy Optimization (GRPO)} GRPO \citep{grpo} bypasses the need for the value model by computing the relative advantage of each response within a group of responses to the same query. Specifically, GRPO optimizes the following objective:

\begin{align}
\label{equ:grpo}
\mathcal{J}_\text{GRPO}(\theta) = \mathbb{E}_{ x \sim \mathcal{D},\, \{y_i\}_{i=1}^G \sim \pi_{\theta_\text{old}}( \cdot | x) }
\left[ \frac{1}{G} \sum_{i=1}^{G} \frac{1}{|y_i|} \sum_{t=1}^{|y_i|} 
\min \left( w_{i,t}(\theta) \widehat{A}_{i,t},  \, \mathrm{clip} \left( w_{i,t}(\theta), 1 - {\varepsilon}, 1 + {\varepsilon}\right) \widehat{A}_{i,t} \right)
\right],
\end{align}

where $G$ is the number of generated responses to each query $x$ (i.e., the group size), and the importance ratio $w_{i,t}(\theta)$ and advantage $\widehat{A}_{i,t}$ of token $y_{i,t}$ are:

\begin{align}
    w_{i,t}(\theta)=\frac{ \pi_{\theta} (y_{i,t} | x, y_{i,<t}) }{ \pi_{\theta_\text{old}} (y_{i,t} | x,y_{i,<t})},\quad\ 
    \widehat{A}_{i,t} = \widehat{A}_{i} = \frac{r(x, y_i) - \mathrm{mean} \left( \{ r(x, y_i) \}_{i=1}^G \right) }{ \mathrm{std} \left( \{ r(x, y_i) \}_{i=1}^G \right) },
\end{align}

respectively, where all the tokens in $y_i$ share the same advantage as $\widehat{A}_{i}$.

\section{Motivation}
\label{sec:motivation}

As model capacities, sparsity patterns (such as those found in Mixture-of-Experts architectures), and generated response lengths have increased, reinforcement learning requires utilizing large rollout batch sizes to fully leverage modern hardware resources. To enhance the efficiency of learning from these samples, it is a common approach to divide the substantial rollout batches into several smaller mini-batches when performing gradient updates. This process inherently leads to an off-policy learning scenario, wherein the generated responses $y$ originate from a prior policy $\pi_{\theta_\text{old}}$ instead of the latest policy $\pi_\theta$ under optimization. Consequently, this shift justifies the implementation of clipping strategies in methods like PPO and GRPO, as these mechanisms restrict the influence of samples that are significantly out-of-distribution relative to the current policy during gradient calculations.

While mechanisms like clipping aim to manage this off-policy discrepancy, we identify a more fundamental issue in GRPO: \textit{its objective is ill-posed}. This problem becomes particularly acute when training large models on long-response tasks, leading to catastrophic model collapse. The ill-posed nature of the GRPO objective stems from a misapplication of importance sampling weights. The principle of importance sampling is to estimate the expectation of a function $f$ under a target distribution $\pi_\text{tar}$ by re-weighting samples drawn from a behavior distribution $\pi_\text{beh}$:

\begin{align}
\mathbb{E}_{z \sim \pi_\text{tar}} \left[ f(z) \right] = \mathbb{E}_{z \sim \pi_\text{beh}} \left[ \frac{ \pi_\text{tar} (z) }{ \pi_\text{beh} (z)} \, f(z) \right].
\end{align}

A key aspect here is that effective correction for the distributional discrepancy depends on averaging over a large number of samples ($N \gg 1$) from the behavior policy $\pi_\text{beh}$, so that the importance weight $\frac{ \pi_\text{tar} (z) }{ \pi_\text{beh} (z)}$ serves its intended function.

By contrast, in GRPO, the importance weight $\frac{ \pi_{\theta} (y_{i,t} | x, y_{i,<t}) }{ \pi_{\theta_\text{old}} (y_{i,t} | x,y_{i,<t})}$ is applied individually at every token position $t$. This approach relies on just one sample $y_{i,t}$ per next-token distribution $\pi_{\theta_\text{old}}( \cdot | x, y_{i, <t})$, thereby failing to sufficiently correct for distribution mismatch. Rather than mitigating bias, this mechanism injects substantial high-variance noise into the gradient estimates, which accumulates over extended sequences and is further aggravated by the presence of the clipping procedure. Through empirical investigations, we have found that this often results in irreversible model collapse. After such collapse, further attempts to train—even reverting to previous checkpoints, carefully adjusting hyperparameters (including clipping bounds), increasing the length of generations, or modifying RL query structures—fail to restore progress.

This finding points to a critical deficiency inherent in the GRPO framework. Specifically, the inadequacy of token-level importance weighting underlines an essential insight: \textbf{the granularity of the optimization objective should align with that of the reward signal}. Because rewards are awarded based on whole sequences, off-policy correction performed at the token level appears fundamentally flawed. This realization leads us to discard the token-level approach in favor of adopting sequence-level importance weights and performing direct optimization at the sequence scale.

\section{Algorithm}

\subsection{GSPO: Group Sequence Policy Optimization}

Although the token-level importance weight $\frac{ \pi_{\theta} (y_{i,t} | x, y_{i,<t}) }{ \pi_{\theta_\text{old}} (y_{i,t} | x,y_{i,<t})}$ presents issues within GRPO, we note that, for language generation tasks, the \textit{sequence-level} importance weight $\frac{ \pi_{\theta} (y | x) }{ \pi_{\theta_\text{old}} (y | x)}$ possesses a well-defined theoretical interpretation: it quantifies the extent to which the response $y$—sampled from $\pi_{\theta_\text{old}} (\cdot | x)$—differs from that generated by $\pi_{\theta} (\cdot | x)$. This aligns with the use of sequence-level rewards and also provides an intuitive foundation for the clipping strategy.

Based on this straightforward observation, we propose the \textbf{Group Sequence Policy Optimization (GSPO)} algorithm. GSPO employs the following sequence-level optimization objective:

\begin{align}
\mathcal{J}_\text{GSPO} (\theta) =
\mathbb{E}_{ x \sim \mathcal{D},\, \{y_i\}_{i=1}^G \sim \pi_{\theta_\text{old}}( \cdot | x) }
\left[ 
\frac{1}{G} \sum_{i=1}^{G}
\min \left( s_{i}(\theta)  \widehat{A}_{i},  \, \mathrm{clip} \left( s_{i}(\theta), 1 - {\varepsilon}, 1 + {\varepsilon} \right) \widehat{A}_{i} \right) 
\right],
\label{equ:gspo}
\end{align}

where we adopt the group-based advantage estimation:

\begin{align}
\widehat{A}_{i} = \frac{r(x, y_i) - \mathrm{mean} \left( \{ r(x, y_i) \}_{i=1}^G \right) }{ \mathrm{std} \left( \{ r(x, y_i) \}_{i=1}^G \right) },
\end{align}

and define the importance ratio $s_{i}(\theta)$ based on sequence likelihood \citep{click}:

\begin{align}
s_{i}(\theta) = \left( \frac{ \pi_{\theta} (y_i | x) }{ \pi_{\theta_\text{old}} (y_i | x)} \right)^{\frac{1}{|y_i|}}
=
\exp \left( \frac{1}{|y_i|} \sum_{t=1}^{|y_i|} \log \frac{ \pi_{\theta} (y_{i,t} | x, y_{i,<t}) }{ \pi_{\theta_\text{old}} (y_{i,t} | x,y_{i,<t})} \right).
\end{align}

Consequently, GSPO implements clipping at the response level rather than at the level of individual tokens, effectively eliminating highly “off-policy” samples from the gradient calculation. This approach aligns more closely with both sequence-level reward assignment and optimization processes. Additionally, length normalization is incorporated into $s_{i}(\theta)$ to help stabilize its variance and keep its values confined within a standard numerical interval. Without this normalization, modifications in the probability of only a few tokens might produce substantial volatility in the importance ratio across entire sequences, resulting in the need for different clipping thresholds for responses of varying lengths. It is also important to observe that, because of differing importance ratio formulations, the scales of clipping ranges used by GSPO and prior methods such as GRPO often vary substantially.

\subsection{Gradient Analysis}
\label{subsec:gradient}

We can derive the gradient of the GSPO objective as follows (clipping is omitted for brevity):

\begin{align}
\nabla_{\theta} \mathcal{J}_\text{GSPO} (\theta)
=&\ 
\nabla_{\theta} \mathbb{E}_{ x \sim \mathcal{D},\, \{y_i\}_{i=1}^G \sim \pi_{\theta_\text{old}}( \cdot | x) }
\left[ 
\frac{1}{G} \sum_{i=1}^{G} s_{i}(\theta) \widehat{A}_{i}
\right] \\
= &\ 
\mathbb{E}_{ x \sim \mathcal{D},\, \{y_i\}_{i=1}^G \sim \pi_{\theta_\text{old}}( \cdot | x) }
\left[ 
\frac{1}{G} \sum_{i=1}^{G}
s_{i}(\theta) \widehat{A}_{i}
\cdot \nabla_{\theta} \log s_{i}(\theta)
\right] \\
=&\ 
\mathbb{E}_{ x \sim \mathcal{D},\, \{y_i\}_{i=1}^G \sim \pi_{\theta_\text{old}}( \cdot | x) }
\left[ 
\frac{1}{G} \sum_{i=1}^{G} 
\left( \frac{ \pi_{\theta} (y_i | x) }{ \pi_{\theta_\text{old}} (y_i | x)} \right)^{\frac{1}{|y_i|}} \widehat{A}_{i} 
\cdot \frac{1}{|y_i|} \sum_{t=1}^{|y_i|} \nabla_{\theta} \log \pi_{\theta} (y_{i,t} | x, y_{i,<t}) 
\right].
\label{equ:gspo_grad}
\end{align}

For comparison, the gradient of the GRPO objective is as follows (note that $\widehat{A}_{i,t} = \widehat{A}_{i}$):

\begin{align}
\nabla_{\theta} \mathcal{J}_\text{GRPO} (\theta) 
=&\ 
\nabla_{\theta} \mathbb{E}_{ x \sim \mathcal{D},\, \{y_i\}_{i=1}^G \sim \pi_{\theta_\text{old}}( \cdot | x) }
\left[ \frac{1}{G} \sum_{i=1}^{G} \frac{1}{|y_i|} \sum_{t=1}^{|y_i|} 
w_{i,t}(\theta) \widehat{A}_{i,t}
\right] \\
=&\
\mathbb{E}_{ x \sim \mathcal{D},\, \{y_i\}_{i=1}^G \sim \pi_{\theta_\text{old}}( \cdot | x) }
\left[ \frac{1}{G} \sum_{i=1}^{G} \widehat{A}_{i} 
\cdot \frac{1}{|y_i|} \sum_{t=1}^{|y_i|} 
\frac{ \pi_{\theta} (y_{i,t} | x, y_{i,<t}) }{ \pi_{\theta_\text{old}} (y_{i,t} | x,y_{i,<t})} 
\nabla_{\theta} \log \pi_{\theta} (y_{i,t} | x, y_{i,<t})  
\right].
\end{align}

The primary difference separating GSPO from GRPO centers on \textit{the method by which each approach assigns weights to the gradients of the token log likelihoods}. GRPO determines each token’s weight based on its corresponding ``importance weight,'' given by $\frac{ \pi_{\theta} (y_{i,t} | x, y_{i,<t}) }{ \pi_{\theta_\text{old}} (y_{i,t} | x,y_{i,<t})}$. Notably, these weights are often disparate—taking values in the intervals $(0, 1+\varepsilon]$ when $\widehat{A}_{i} > 0$ or $[1-\varepsilon, +\infty)$ when $\widehat{A}_{i} < 0$—and thus should not be regarded as insignificant. Over time, as training unfolds, the effects of these uneven weights can accumulate, introducing potential instability or unpredictable learning dynamics. Conversely, GSPO uniformly weights every token within a generated response, thereby avoiding the instability introduced by the token-level importance weighting present in GRPO.

\subsection{GSPO-token: A Token-level Objective Variant}

In scenarios like multi-turn RL, we may desire a finer-grained advantage adjustment than the sequence level. To this end, we introduce a token-level objective variant of GSPO, namely \textbf{GSPO-token}, to allow token-wise advantage customization:

\begin{align}
\mathcal{J}_\text{GSPO-token}(\theta)
=
\mathbb{E}_{ x \sim \mathcal{D},\, \{y_i\}_{i=1}^G \sim \pi_{\theta_\text{old}}( \cdot | x) }
\left[ 
\frac{1}{G} \sum_{i=1}^{G} \frac{1}{|y_i|} \sum_{t=1}^{|y_i|} 
\min \left( s_{i,t}(\theta) \widehat{A}_{i,t},  \, \mathrm{clip} \left( s_{i,t}(\theta), 1 - {\varepsilon}, 1 + {\varepsilon} \right) \widehat{A}_{i,t} \right) 
\right],
\label{equ:gspo-token}
\end{align}

where

\begin{align}
s_{i,t}(\theta) 
= \mathrm{sg} \left[ s_{i}(\theta) \right]  \cdot \frac{ \pi_{\theta} (y_{i,t} | x, y_{i,<t}) }{ \mathrm{sg} \left[ \pi_{\theta} (y_{i,t} | x, y_{i,<t}) \right] },
\end{align}

and $\mathrm{sg}[\cdot]$ denotes only taking the numerical value but stopping the gradient, corresponding to the \texttt{detach} operation in PyTorch. The gradient of GSPO-token can be derived as:

\begin{align}
\nabla_{\theta} \mathcal{J}_\text{GSPO-token}(\theta)
=&\ 
\nabla_{\theta} \mathbb{E}_{ x \sim \mathcal{D},\, \{y_i\}_{i=1}^G \sim \pi_{\theta_\text{old}}( \cdot | x) }
\left[ 
\frac{1}{G} \sum_{i=1}^{G}
\frac{1}{|y_i|} \sum_{t=1}^{|y_i|} 
s_{i,t}(\theta) \widehat{A}_{i,t}
\right] \\
=&\ 
\mathbb{E}_{ x \sim \mathcal{D},\, \{y_i\}_{i=1}^G \sim \pi_{\theta_\text{old}}( \cdot | x) }
\left[ 
\frac{1}{G} \sum_{i=1}^{G} s_i (\theta) \cdot \frac{1}{|y_i|} \sum_{t=1}^{|y_i|} 
\widehat{A}_{i,t} \frac{ \nabla_{\theta} \pi_{\theta} (y_{i,t} | x, y_{i,<t}) }{ \pi_{\theta} (y_{i,t} | x, y_{i,<t}) }
\right] \\
=&\ 
\mathbb{E}_{ x \sim \mathcal{D},\, \{y_i\}_{i=1}^G \sim \pi_{\theta_\text{old}}( \cdot | x) }
\left[ 
\frac{1}{G} \sum_{i=1}^{G}  \left( \frac{ \pi_{\theta} (y_i | x) }{ \pi_{\theta_\text{old}} (y_i | x)} \right)^{\frac{1}{|y_i|}}
\cdot \frac{1}{|y_i|} \sum_{t=1}^{|y_i|}
\widehat{A}_{i,t} \nabla_{\theta} \log \pi_{\theta} (y_{i,t} | x, y_{i,<t})  
\right].
\label{equ:gspo-token_grad}
\end{align}

Observe that $\frac{ \pi_{\theta} (y_{i,t} | x, y_{i,<t}) }{ \mathrm{sg} \left[ \pi_{\theta} (y_{i,t} | x, y_{i,<t}) \right] }$ evaluates to 1, which means that $s_{i,t}(\theta)$ is, in practice, equivalent to $s_{i}(\theta)$ in value. A comparison between Equation~\eqref{equ:gspo} and~\eqref{equ:gspo-token}, as well as Equation~\eqref{equ:gspo_grad} and~\eqref{equ:gspo-token_grad}, indicates that GSPO-token and GSPO yield the same results regarding their optimization objective, the clipping condition, and the formal gradient expression, as long as all tokens in the response $y_i$ are assigned an identical advantage (specifically, $\widehat{A}_{i,t} = \widehat{A}_{i}$ for each $t$). However, GSPO-token provides greater adaptability by permitting distinct advantage values for each token.

\section{Experiments and Discussion}

\subsection{Empirical Results}

Our experimental setup employs a cold-start model that is fine-tuned from Qwen3-30B-A3B-Base, and we present both the training reward trajectories and the evolution of model performance across several benchmarks: AIME'24 (mean Pass@1 over 32 trials), LiveCodeBench (202410-202502, mean Pass@1 across 8 samples), and CodeForces (Elo Rating). Throughout RL training, rollout data in each batch is divided into four distinct mini-batches to perform gradient updates. For GSPO, the left and right clip values in Equation~\eqref{equ:gspo} are chosen as 3e-4 and 4e-4, respectively. For a comparative baseline, we employ GRPO, assigning its left and right clipping parameters in Equation~\eqref{equ:grpo} to 0.2 and 0.27, respectively, with hyperparameters optimized for equitable assessment. It is important to highlight that GRPO requires implementation of the Routing Replay training approach in order for MoE RL to converge correctly; this aspect is elaborated in \S~\ref{subsec:routing_replay}. In contrast, \textit{GSPO eliminates this requirement}.

As depicted in Figure~\ref{fig:results}, the training process using GSPO exhibits consistent stability. Our findings indicate that \textbf{GSPO enables ongoing improvements in performance by scaling up training compute, periodically refreshing the query set, and lengthening the generated sequences}. In addition, GSPO outperforms GRPO with respect to training efficiency, attaining higher training accuracy and superior results on benchmarks while utilizing equivalent training compute and an equal number of consumed queries. Lastly, we have effectively deployed GSPO in the RL training of the newest Qwen3 models, thereby providing strong evidence that GSPO successfully harnesses the potential of RL scaling for large language models.

\subsection{Curious Observation on Clipping Fractions}

An important difference between GSPO and GRPO lies in their approaches to clipping: GSPO applies clipping to whole responses, whereas GRPO does so at the level of individual tokens. As depicted in Figure~\ref{fig:clipping}, GSPO results in a fraction of clipped tokens that is roughly two orders of magnitude greater than that of GRPO; this substantial gap in clipping rate persists regardless of any adjustments made to the clipping ranges. Nevertheless, even though GSPO discards a much greater proportion of tokens—resulting in fewer tokens being available for training or gradient calculation—it still outperforms GRPO in terms of training efficiency. This seemingly paradoxical observation—where significantly increased token clipping correlates with enhanced training efficiency—suggests that GRPO’s gradient estimates at the token level are fundamentally noisy and suboptimal for effective exploitation of sampled data. Conversely, by operating at the level of entire sequences, GSPO generates a learning signal that is both more robust and conducive to efficient training.

\subsection{Benefit of GSPO for MoE Training}
\label{subsec:routing_replay}

\paragraph{Background}
MoE models, due to their inherently sparse activation patterns, pose distinct stability issues in RL training compared to their dense model counterparts. Specifically, our observations reveal that employing the GRPO algorithm with MoE architectures leads to pronounced \textit{expert-activation volatility}, which can undermine the successful convergence of RL procedures. More precisely, post one or more rounds of gradient updates, there can be substantial shifts in which experts are selected for generating identical responses. As an illustration, consider the 48-layer Qwen3-30B-A3B-Base model: following each RL-driven gradient update, approximately 10\% of the activated experts for a given rollout sample under the updated policy $\pi_{\theta}$ differ from those under the previous policy $\pi_{\theta_\text{old}}$. This issue is further amplified in deeper MoE architectures, resulting in highly unstable token-level importance ratios $w_{i,t}(\theta) = \frac{ \pi_{\theta} (y_{i,t} | x, y_{i,<t}) }{ \pi_{\theta_\text{old}} (y_{i,t} | x,y_{i,<t})}$ and, as elaborated in \S~\ref{sec:motivation} and~\ref{subsec:gradient}, renders these ratios unreliable—ultimately disrupting regular RL training convergence.

\paragraph{Our Previous Approach} In our earlier work addressing this problem, we utilized a training method called \textbf{Routing Replay}. The method operates by saving the set of experts activated under $\pi_{\theta_\text{old}}$ and reapplying these same expert selections—effectively ``replaying'' the routing—within $\pi_{\theta}$ as we compute the importance weights $w_{i,t}(\theta) = \frac{ \pi_{\theta} (y_{i,t} | x, y_{i,<t}) }{ \pi_{\theta_\text{old}} (y_{i,t} | x,y_{i,<t})}$. By synchronizing the choice of activated network for each token $y_{i,t}$ in both $\pi_{\theta} (y_{i,t} | x, y_{i,<t})$ and $\pi_{\theta_\text{old}} (y_{i,t} | x,y_{i,<t})$, we maintain consistency in the computed token-level importance ratios and are able to ensure stable optimization for the recurring network configuration throughout the gradient updates. As depicted in Figure~\ref{fig:routing_replay}, this Routing Replay method is pivotal to achieving stable convergence in the GRPO training process for MoE architectures.

\paragraph{Benefit of GSPO} While Routing Replay ensures successful convergence during GRPO training of MoE models, it does so by reusing routing decisions, which leads to increased memory consumption and communication costs, as well as restricting the model’s true capacity. In contrast, as depicted in Figure~\ref{fig:results}, GSPO removes the reliance on Routing Replay altogether. It can straightforwardly calculate the importance ratios $s_i(\theta)$ as usual, maintain convergence, and achieve stable optimization. The fundamental reason is that GSPO exclusively considers sequence-level likelihoods (i.e., $\pi_{\theta} (y_i | x)$) rather than individual token-level probabilities (i.e., $\pi_{\theta} (y_{i,t} | x, y_{i,<t})$). Given that the MoE architecture retains its language modeling proficiency, the sequence likelihood tends to remain steady and does not exhibit significant volatility. Ultimately, GSPO directly addresses the instability arising from expert activation in MoE models, eliminating the necessity for elaborate solutions such as Routing Replay. This approach both streamlines and stabilizes training and enables the MoE model to fully utilize its potential without imposed limitations.

\subsection{Benefit of GSPO for RL Infrastructure}

Due to precision mismatches that often arise between training frameworks (such as Megatron) and inference frameworks (such as SGLang and vLLM), standard practice usually involves recalculating the likelihoods of sampled outputs under the prior policy $\pi_{\theta_\text{old}}$ using the training engine. Nonetheless, GSPO operates at the level of complete sequences rather than individual tokens, and sequence-level likelihoods are generally far less sensitive to small precision inconsistencies. As a result, GSPO permits the direct utilization of likelihoods produced by inference engines during optimization, eliminating the necessity of recomputation with the training engine. This feature can provide significant advantages in use cases like partial rollout, multi-turn RL, and within frameworks where training and inference are separated.

\section{Conclusion}

We introduce Group Sequence Policy Optimization (GSPO), an innovative reinforcement learning method designed to train large language models. Built upon the core ideas of importance sampling, GSPO constructs importance ratios using sequence likelihood, implementing clipping, rewarding, and optimization at the sequence level. This algorithm offers significant advantages in terms of training stability, sample efficiency, and overall effectiveness relative to GRPO, and proves especially advantageous in the large-scale reinforcement learning of MoE models. This forms the basis for the remarkable progress observed in the newest Qwen3 models. With GSPO serving as a robust, scalable framework, we aim to further expand RL methodologies and anticipate groundbreaking developments in machine intelligence.

\end{document}